\documentclass[12pt]{article}
\usepackage[utf8]{inputenc}
\usepackage[english]{babel}
\usepackage{amsmath}
\usepackage{graphicx}
\usepackage{hyperref}
\usepackage{listings}
\usepackage{xcolor}
\usepackage{booktabs}
\usepackage{algorithm}
\usepackage{algpseudocode}

% Code listing style
\lstset{
    basicstyle=\ttfamily\footnotesize,
    breaklines=true,
    frame=single,
    numbers=left,
    numberstyle=\tiny\color{gray},
    captionpos=b
}

\title{LLM-Based Test Smell Refactoring Methodology}
\author{}
\date{}

\begin{document}

\maketitle

\section{LLM-Based Refactoring Methodology}
\label{sec:llm-methodology}

\subsection{Overview}

This research employs Large Language Models (LLMs) as automated refactoring agents to systematically eliminate test smells from JavaScript codebases. The methodology integrates state-of-the-art code-specialized LLMs with structured prompt engineering techniques to transform smelly test code into maintainable, best-practice-compliant implementations.

\subsection{Motivation and Research Context}

Traditional test smell refactoring relies heavily on manual intervention, requiring developers to identify patterns, understand antipatterns, and apply contextual fixes~\cite{mockus2022test}. This process is time-consuming, error-prone, and does not scale to large codebases. Recent advances in Large Language Models, particularly those fine-tuned for code generation~\cite{brown2020language, chen2021evaluating}, present an opportunity to automate this refactoring process while maintaining semantic equivalence and test validity.

Our approach addresses three key challenges:
\begin{enumerate}
    \item \textbf{Semantic Preservation}: Ensuring refactored tests maintain identical behavioral specifications
    \item \textbf{Smell Elimination}: Systematically removing detected antipatterns without introducing new smells
    \item \textbf{Generalization}: Creating a methodology applicable across diverse test smell categories
\end{enumerate}

\subsection{LLM Selection and Rationale}

The research employed six distinct LLMs across multiple providers, selected based on code generation capabilities, context window capacity, and architectural diversity (Table~\ref{tab:llm-models}).

\begin{table}[h]
\centering
\caption{LLM Models Employed in Refactoring Experiments}
\label{tab:llm-models}
\begin{tabular}{@{}llll@{}}
\toprule
\textbf{Model} & \textbf{Provider} & \textbf{Parameters} & \textbf{Specialization} \\ \midrule
Qwen 2.5 Coder 32B & HuggingFace & 32B & Code generation (primary) \\
CodeLlama 34B & Meta/HF & 34B & Code completion \\
Claude Sonnet 4.6 & Anthropic & N/A & Reasoning \& code \\
GPT-5.2 & OpenAI & N/A & General-purpose \\
DeepSeek-Coder 33B & DeepSeek & 33B & Code-specialized \\
Gemini 1.5 Pro & Google & N/A & Long-context reasoning \\ \bottomrule
\end{tabular}
\end{table}

\textbf{Qwen 2.5 Coder 32B Instruct} was designated as the primary model due to its superior performance on code refactoring benchmarks, extensive training on GitHub repositories, and optimized instruction-following capabilities for code transformation tasks~\cite{zheng2024qwen2}.

\subsection{Prompt Engineering Strategies}
\label{sec:prompt-strategies}

The research implements three distinct prompt engineering techniques, each representing a progressive sophistication in guiding LLM reasoning processes:

\subsubsection{Zero-Shot Prompting}

Zero-shot prompting provides the LLM with a direct instruction and smell definition without demonstrative examples. This baseline strategy evaluates the model's inherent understanding of refactoring principles.

\textbf{Structure:}
\begin{itemize}
    \item \textbf{System Message}: Role definition as ``JavaScript test refactoring expert''
    \item \textbf{Smell Definition}: Academic definition from literature
    \item \textbf{Original Code}: Smelly test snippet
    \item \textbf{Output Constraints}: Strict format requirements (JavaScript code block only)
\end{itemize}

\textbf{Mathematical Formulation:}
\begin{equation}
    P_{\text{zero}}(C_r | C_s, D_s) = \text{LLM}(\text{Task} \oplus D_s \oplus C_s)
\end{equation}
where $C_r$ is refactored code, $C_s$ is smelly code, $D_s$ is smell definition, and $\oplus$ denotes prompt concatenation.

\textbf{Advantages:} Minimal prompt complexity, evaluates intrinsic model knowledge, computationally efficient.

\textbf{Limitations:} May fail to capture domain-specific refactoring patterns not seen during pre-training.

\subsubsection{Few-Shot Prompting}

Few-shot prompting augments the zero-shot approach with 1--3 curated examples demonstrating smell-to-clean transformations. This technique leverages in-context learning to guide the model toward expected refactoring patterns.

\textbf{Structure:}
\begin{itemize}
    \item \textbf{Examples Section}: Original (smelly) $\rightarrow$ Refactored (clean) pairs
    \item \textbf{Example Selection}: Maximum 3 examples from smell catalog
    \item \textbf{Pattern Emphasis}: ``Study the examples to understand the refactoring pattern''
\end{itemize}

\textbf{Mathematical Formulation:}
\begin{equation}
    P_{\text{few}}(C_r | C_s, D_s, E) = \text{LLM}(\text{Task} \oplus D_s \oplus E_{1:k} \oplus C_s)
\end{equation}
where $E_{1:k}$ represents $k$ example pairs ($k \leq 3$).

\textbf{Example Structure:}
\begin{lstlisting}[language=,caption=Few-Shot Example Format]
### Example 1
Original (with Anonymous Test):
```javascript
it('should handle date', () => { /* ... */ });
```

Refactored (smell removed):
```javascript
it('formats ISO date string to month day, year format', () => {
    /* ... */
});
```
\end{lstlisting}

\textbf{Advantages:} Improves consistency, reduces ambiguity in refactoring intent, enables pattern transfer.

\textbf{Limitations:} Example selection bias, increased prompt token cost.

\subsubsection{Chain-of-Thought (CoT) Prompting}

Chain-of-Thought prompting implements a structured reasoning framework where the LLM is instructed to internally decompose the refactoring task into explicit cognitive steps before generating the final output. This approach is inspired by \cite{wei2022chain} and adapted for code refactoring.

\textbf{Reasoning Steps (Hidden):}
\begin{enumerate}
    \item \textbf{Locate the Smell}: Identify manifestation based on detection criteria
    \item \textbf{Understand Intent}: Determine test behavior and assertion semantics
    \item \textbf{Evaluate Impact}: Assess why structure violates best practices
    \item \textbf{Plan Refactoring}: Design solution using provided strategies
    \item \textbf{Validate}: Ensure requirements satisfaction
\end{enumerate}

\textbf{Additional Components:}
\begin{itemize}
    \item \textbf{Detection Criteria}: Formal smell identification rules
    \item \textbf{Refactoring Strategies}: Operational guidance catalog
    \item \textbf{Optional Examples}: Maximum 2 reference examples
\end{itemize}

\textbf{Mathematical Formulation:}
\begin{equation}
    P_{\text{cot}}(C_r | C_s, D_s, R_s, S_d, E) = \text{LLM}(\text{Task} \oplus D_s \oplus S_d \oplus R_s \oplus E_{1:2} \oplus C_s \oplus \text{CoT}_{\text{template}})
\end{equation}
where $R_s$ denotes refactoring strategies, $S_d$ represents detection criteria, and $\text{CoT}_{\text{template}}$ is the step-by-step reasoning framework.

\textbf{Prompt Fragment:}
\begin{lstlisting}[caption=CoT Internal Reasoning Template]
### Internal Reasoning (Hidden Chain-of-Thought)

You MUST internally reason step by step:

1. **Locate the Smell**: Identify manifestation based on detection
2. **Understand Intent**: Determine test behavior
3. **Evaluate Impact**: Assess why structure violates practices
4. **Plan Refactoring**: Design solution using strategies
5. **Validate**: Ensure requirements met

**CRITICAL:** DO NOT reveal reasoning. ONLY output final code.
\end{lstlisting}

\textbf{Advantages:} Enhanced reasoning transparency, improved handling of complex smells, systematic approach enforcement.

\textbf{Comparison of Strategies:}
Table~\ref{tab:prompt-comparison} contrasts the three approaches across key dimensions.

\begin{table}[h]
\centering
\caption{Comparative Analysis of Prompt Engineering Strategies}
\label{tab:prompt-comparison}
\begin{tabular}{@{}llll@{}}
\toprule
\textbf{Dimension} & \textbf{Zero-Shot} & \textbf{Few-Shot} & \textbf{Chain-of-Thought} \\ \midrule
Prompt Tokens & Low ($\sim$200) & Medium ($\sim$600) & High ($\sim$800) \\
Examples & 0 & 1--3 & 0--2 \\
Reasoning & Implicit & Implicit & Explicit (hidden) \\
Complexity & Minimal & Moderate & High \\
Guidance & Definition only & Patterns & Strategies + Structure \\
Consistency & Variable & Improved & Highest \\ \bottomrule
\end{tabular}
\end{table}

\subsection{Output Constraints and Validation}

All prompting strategies enforce strict output requirements to ensure extractability and executability:

\begin{itemize}
    \item \textbf{Format}: JavaScript code block only (markdown fenced: \texttt{```javascript})
    \item \textbf{Completeness}: Full test function (entire \texttt{it()} or \texttt{describe()} block)
    \item \textbf{Semantic Preservation}: Identical assertions and behavioral specifications
    \item \textbf{Syntactic Correctness}: Executable code without syntax errors
    \item \textbf{Prohibited Content}: No explanations, comments, or descriptive text outside code
\end{itemize}

\subsection{Automated Experiment Pipeline}
\label{sec:pipeline}

The refactoring experiments execute through a two-phase automated pipeline designed for reproducibility and scalability:

\subsubsection{Phase 1: Refactoring (LLM Generation)}

\begin{algorithm}
\caption{Batch Refactoring Phase}
\label{alg:refactor-phase}
\begin{algorithmic}[1]
\Require Smell dataset $\mathcal{S} = \{s_1, \ldots, s_n\}$, strategy $\sigma \in \{1, 2, 3\}$, model $M$
\Ensure Experiment records with refactored code
\For{$s_i \in \mathcal{S}$}
    \State $context \gets \text{LoadSmellContext}(s_i)$ \Comment{Definition, examples, strategies}
    \State $prompt \gets \text{BuildPrompt}(\sigma, context)$ \Comment{Apply strategy}
    \State $response \gets M.\text{generate}(prompt)$ \Comment{LLM API call}
    \State $code \gets \text{ExtractCode}(response)$ \Comment{Regex + fence parsing}
    \State $exp \gets \text{CreateExperiment}(s_i, \sigma, M, code)$
    \State $\text{SaveToDatabase}(exp)$ \Comment{Marks \texttt{refactor\_phase\_completed}}
\EndFor
\State $manifest \gets \text{GenerateManifest}(\mathcal{E})$ \Comment{JSON file for Phase 2}
\end{algorithmic}
\end{algorithm}

\textbf{Key Operations:}
\begin{itemize}
    \item \textbf{Context Loading}: Retrieves smell definition, detection criteria, refactoring strategies, and examples from catalog
    \item \textbf{Prompt Construction}: Dynamically builds system + user messages based on strategy
    \item \textbf{Code Extraction}: Applies regex patterns to isolate JavaScript code from markdown fences
    \item \textbf{Database Persistence}: Stores original code, refactored code, LLM metrics (tokens, latency)
    \item \textbf{Manifest Generation}: Creates JSON file mapping experiment IDs for Phase 2 execution
\end{itemize}

\subsubsection{Phase 2: Validation and Analysis}

\begin{algorithm}
\caption{Batch Execution and Validation Phase}
\label{alg:execute-phase}
\begin{algorithmic}[1]
\Require Experiment set $\mathcal{E}$ from manifest or database
\Ensure Validated experiments with test results and smell detection
\For{$e \in \mathcal{E}$}
    \State $\text{CreateBackup}(e.\text{repository}, e.\text{file\_path})$
    \State $\text{ApplyRefactoring}(e.\text{refactored\_code}, e.\text{file\_path})$
    \State $smell\_results \gets \text{RunSmellDetection}(e.\text{repository})$
    \State $test\_results \gets \text{RunTestSuite}(e.\text{repository})$
    \State $analysis \gets \text{CompareResults}(baseline, smell\_results, test\_results)$
    \State $\text{SaveResults}(e, analysis)$ \Comment{Marks \texttt{execution\_phase\_completed}}
    \State $\text{RestoreBackup}(e.\text{repository}, e.\text{file\_path})$
\EndFor
\end{algorithmic}
\end{algorithm}

\textbf{Validation Metrics:}
\begin{enumerate}
    \item \textbf{Test Pass Rate}: $\frac{\text{tests\_passed}}{\text{tests\_total}}$ (baseline vs. refactored)
    \item \textbf{Code Coverage}: Statement, branch, function, line coverage percentages
    \item \textbf{Smell Removal}: Binary flag indicating target smell elimination
    \item \textbf{New Smell Introduction}: Detected new smells not present in baseline
    \item \textbf{Test Execution Time}: Performance impact measurement
\end{enumerate}

\subsection{Database Schema and Traceability}

Each experiment generates a comprehensive record ensuring full traceability:

\begin{table}[h]
\centering
\caption{Experiment Database Schema (Key Fields)}
\label{tab:schema}
\begin{tabular}{@{}ll@{}}
\toprule
\textbf{Field} & \textbf{Description} \\ \midrule
\texttt{study\_smell\_id} & Reference to detected smell \\
\texttt{prompting\_approach} & Strategy name (Zero-Shot/Few-Shot/CoT) \\
\texttt{ai\_model\_version} & LLM identifier \\
\texttt{original\_code} & Baseline smelly code \\
\texttt{refactored\_code} & LLM-generated code \\
\texttt{smell\_removed} & Boolean: target smell eliminated \\
\texttt{introduced\_new\_smells} & Boolean: new smells detected \\
\texttt{tests\_still\_passing} & Boolean: test suite status \\
\texttt{coverage\_decreased} & Boolean: coverage regression \\
\texttt{tokens\_used} & Total prompt + completion tokens \\
\texttt{llm\_latency\_seconds} & API response time \\ \bottomrule
\end{tabular}
\end{table}

\subsection{Statistical Experimental Design}

The research follows a factorial design evaluating:
\begin{itemize}
    \item \textbf{Independent Variables}: 
        \begin{itemize}
            \item Prompt strategy ($\sigma \in \{1, 2, 3\}$)
            \item LLM model ($M \in \{\text{6 models}\}$)
            \item Smell type ($T \in \{\text{10 smell categories}\}$)
        \end{itemize}
    \item \textbf{Dependent Variables}:
        \begin{itemize}
            \item Smell removal success rate
            \item Test preservation rate
            \item Coverage maintenance rate
            \item New smell introduction rate
        \end{itemize}
\end{itemize}

\textbf{Experiment Space:}
\begin{equation}
    |\mathcal{E}| = |\mathcal{S}| \times |\Sigma| \times |M| = n \times 3 \times 6
\end{equation}
where $n$ is the number of study smells selected for the corpus.

\subsection{Reproducibility Mechanisms}

To ensure replicability:
\begin{enumerate}
    \item \textbf{Version Control}: All refactored code and experiment metadata stored in Git
    \item \textbf{API Determinism}: Fixed temperature (0.1) and top-p (0.95) parameters
    \item \textbf{Manifest Files}: JSON manifests link experiments to exact LLM responses
    \item \textbf{Backup System}: Automatic file restoration ensures repository integrity
    \item \textbf{Database Export}: Complete experiment database exportable to SQLite/CSV
\end{enumerate}

\subsection{Ethical Considerations}

The research adheres to open-source software analysis guidelines:
\begin{itemize}
    \item All analyzed repositories are publicly available under permissive licenses
    \item No proprietary or confidential code is included in the dataset
    \item LLM providers are used in accordance with their terms of service
    \item Generated refactorings are not automatically committed to upstream projects
\end{itemize}

\subsection{Limitations and Threats to Validity}

\subsubsection{Internal Validity}
\begin{itemize}
    \item \textbf{LLM Non-Determinism}: Despite fixed parameters, LLMs may produce slight variations across runs
    \item \textbf{Smell Detector Accuracy}: False positives/negatives in baseline smell detection propagate to experiments
\end{itemize}

\subsubsection{External Validity}
\begin{itemize}
    \item \textbf{Generalization}: Results specific to JavaScript/Jest ecosystem; transferability to other languages/frameworks requires validation
    \item \textbf{Repository Selection Bias}: Dataset limited to popular GitHub projects with high test coverage
\end{itemize}

\subsubsection{Construct Validity}
\begin{itemize}
    \item \textbf{Smell Definition Subjectivity}: Academic definitions may not capture all real-world smell manifestations
    \item \textbf{Human Evaluation}: Manual inspection of subset needed to validate automated metrics
\end{itemize}

\subsection{Summary}

This methodology establishes a rigorous, reproducible framework for evaluating LLM capabilities in automated test smell refactoring. By systematically varying prompt engineering strategies and LLM architectures across a diverse smell corpus, the research quantifies:
\begin{enumerate}
    \item The comparative effectiveness of zero-shot, few-shot, and chain-of-thought prompting for code refactoring tasks
    \item The relative performance of six state-of-the-art LLMs on smell elimination
    \item The preservation of test semantics and coverage under automated refactoring
\end{enumerate}

The two-phase pipeline, comprehensive database schema, and validation metrics provide a foundation for future work in LLM-assisted software maintenance and test quality improvement.

\end{document}
